\documentclass{../tex/ficheseance}
\usepackage[french]{babel}

\numfiche{02B}
\datefiche{17/10/2026}
\archer{Groupe compétition}
\theme{Régularité du placement du coude d'armement}
\objectifseance{Stabiliser le geste sous pression à 18m}

\begin{document}

	\makeficheheader

	\begin{tableauSeance}

		\ligneSeance
		{\textbf{Échauffement \& Prise de repères} : Mobiliser la ceinture scapulaire et ressentir le placement du coude.} % 1. Objectif
		{15 min} % 2. Temps
		{3 volées à 18m sur blason de 80 cm.} % 3. Situations
		{Élastique, miroir.}         % 4. Outils
		{Mettre l'accent sur le placement du coude d'armement à la paille.}              % 5. Consignes
		{80\% de flèches dans la zone jaune.} % 6. Critères de réussite
		{Stabilité de la tête à l'échappée.}  % 7. Évaluations}

		\ligneSeance
		{\textbf{Regard fixe sur la cible}}
		{15 min}
		{3 volées à 18m sur blason de 80 cm}
		{Élastique, miroir}
		{\begin{itemize}
				\item Mettre l'accent sur la paille
				\item Libération fluide
		\end{itemize}}              % 5. Consignes (avec liste à puces LaTeX !)
		{80\% de flèches dans la zone jaune}
		{Stabilité de la tête à l'échappée}

		\ligneSeance
		{\textbf{Fixation technique :} Conserver la hauteur et l'alignement du coude pendant la phase de maintien/ancrage.}
		{30 min}
		{6 volées de 3 flèches à 18m sur blasons de 40 cm}
		{Caméra / Téléphone (vidéo retardée)}
		{Marquer une pause de 2s à l'ancrage.}
		{Alignement coude/flèche conservé sur au moins 15 flèches sur 18}
		{Absence de variation de la hauteur du coude.}

		\ligneSeance
		{\textbf{Transfert \& Pression :} Valider la stabilité du placement sous contrainte temporelle.}
		{15 min}
		{3 volées de 3 flèches comptées en situation de duel}
		{Feuille de marque / Chronomètre}
		{Appliquer la même routine sous contrainte de temps (20s par flèche)}
		{Score supérieur ou égal à 8 points de moyenne par volée}
		{Maintien du rythme de tir sans précipitation sur la libération.}

	\end{tableauSeance}

\end{document}